% Generated draft: requires mathematical and source review.
\documentclass[11pt]{article}
\usepackage[margin=25mm]{geometry}
\usepackage{fontspec}
\setmainfont{Latin Modern Roman}
\usepackage{amsmath,amssymb,mathtools,unicode-math}
\setmathfont{Latin Modern Math}
\usepackage{enumitem,needspace,xurl,hyperref}
\setlength{\parindent}{0pt}
\setlength{\parskip}{5pt}
\setlength{\emergencystretch}{3em}
\newcommand{\N}{\mathbb{N}}
\newcommand{\Z}{\mathbb{Z}}
\newcommand{\Q}{\mathbb{Q}}
\newcommand{\R}{\mathbb{R}}
\newcommand{\C}{\mathbb{C}}
\newcommand{\F}{\mathbb{F}}
\newcommand{\A}{\mathbb{A}}
\newcommand{\PP}{\mathbb{P}}
\newcommand{\E}{\mathbb{E}}
\newcommand{\eps}{\varepsilon}
\newcommand{\dd}{\,\mathrm{d}}
\begin{document}
\section*{UnsolvedMath ID: 30001026}
\textbf{Source problem code:} catalogue entry
\subsection{Definitions and mathematical statement}
Let \(M\) be a compact Riemannian surface, let \(N_{x,y}(\lambda)\) denote its spectral function, and let \(\Gamma_{x,y}\) be the geodesic segments from \(x\) to \(y\). For a unit-speed segment \(\gamma\), write \(l(\gamma)\) for its length and \(m(\gamma)\) for its Morse index with the intended endpoint convention. The orthogonal Jacobi field in the source has scalar normal component \(J\) normalized by \(J(0)=0\) and \(J'(0)=1\).

For finite \(p\geq1\), the usual Besicovitch seminorm on functions of \(\lambda\geq0\) is
\[
\|h\|_{B^p}=\left(\limsup_{T\to\infty}\frac1T\int_0^T|h(\lambda)|^p\,\dd\lambda\right)^{1/p}.
\]
The Besicovitch almost-periodic class is obtained by approximation by trigonometric polynomials in this seminorm, identifying functions of seminorm zero; precise completion conventions should be checked against the source.

The question is whether there exists \(p\geq1\), chosen for the surface before the points, such that for every \(x\in M\) and Riemannian-volume-almost every \(y\in M\),
\[
\lambda\longmapsto\frac{N_{x,y}(\lambda)}{\sqrt{1+\lambda}}
\]
has the stated geodesic-sum asymptotic expansion in \(B^p\), with frequencies \(l(\gamma)\), phases determined by \(m(\gamma)\), and amplitudes determined by \(l(\gamma)\) and the normalized orthogonal Jacobi field, for \(\gamma\in\Gamma_{x,y}\). The exceptional set of \(y\) may depend on \(x\); the source does not require one common full-measure set for all \(x\).

The extraction gives no actual coefficient, phase, or summation formula. It also omits whether \(\lambda\) cuts off Laplace eigenvalues or their square roots, the spectral-projector normalization, and the precise meaning of the asymptotic expansion in \(B^p\). These omissions prevent writing an exact geodesic series without inventing data.
\subsection{Short English statement}
Seek one Besicovitch exponent for a compact surface so that every fixed basepoint has the proposed off-diagonal geodesic expansion for almost every endpoint. The geometric ingredients and quantifiers are preserved, but the series coefficients and spectral conventions are missing.
\subsection{Sources}
\begin{itemize}
\item Local UnsolvedMath registry, permanent ID \(30001026\), source code \texttt{catalogue entry}, ``Geodesic Expansions of Off-Diagonal Spectral Functions.'' Registry curation: Ulam AI / UnsolvedMath Contributors, CC BY 4.0; underlying sources retain their supplied provenance.
\item \emph{Geometric Group Theory, Hyperbolic Dynamics and Symplectic Geometry} (2009), Oberwolfach Reports, pp.\ 2337--2406, DOI \url{https://doi.org/10.4171/owr/2008/42}.
\item Lapointe, Polterovich and Safarov, supplied title \emph{average growth of the spectral function}: \url{https://arxiv.org/abs/0803.4171}.
\item Supplied literature triage, \texttt{literature\_checked\_at}: 2026-08-21: labeled partially solved; says the Besicovitch almost-periodic expansion is proved for round spheres and flat tori, while existence of a suitable exponent for general surfaces and manifolds remains open. This is supplied dated curation, not independently verified; the statement here retains the surface target rather than introducing an unspecified higher-dimensional normalization.
\end{itemize}
\end{document}
